% ------------------------------------------------------------------------------
% Metodologia
% ------------------------------------------------------------------------------
\chapter{Metodologia}
\label{chap_metodologia}
Este capítulo apresenta os procedimentos adotados para a realização dos experimentos computacionais de queda livre, detalhando os materiais utilizados, os parâmetros de entrada, o método de simulação e a análise dos dados de acurácia.

\section{Materiais Computacionais e Organização dos Dados}
As simulações foram realizadas em ambiente Python, utilizando as bibliotecas \texttt{numpy} para cálculos numéricos, \texttt{pandas} para manipulação de dados e \texttt{matplotlib} para geração de gráficos. Todo o código foi estruturado em um notebook Jupyter, permitindo a execução sequencial dos experimentos e a documentação dos resultados. O código-fonte completo está disponível em \cite{codigoqueda2026}.

Os dados de entrada das simulações consistem nos parâmetros iniciais de cada experimento, como altura inicial ($H$), velocidade inicial ($V$), aceleração da gravidade ($g$), passo de tempo ($\Delta t$), coeficiente de resistência do ar ($c$) e demais variáveis relevantes para cada cenário (queda livre simples, corpo rígido e resistência do ar linear e quadrática).
\section{Procedimento Experimental}


Os experimentos computacionais foram realizados em três etapas sequenciais, cada uma abordando um aspecto físico distinto do movimento:
\begin{enumerate}
	\item \textbf{Queda livre simples:} Simulação do movimento de uma partícula em queda livre, desconsiderando qualquer tipo de resistência do ar. Esta etapa serve como base para validação do método numérico e comparação com a solução analítica clássica.
	\item \textbf{Corpo rígido:} Análise do movimento considerando um corpo rígido, discutindo as diferenças conceituais e possíveis impactos na modelagem, ainda sem resistência do ar.
	\item \textbf{Resistência do ar:} Inclusão da resistência do ar no modelo, tanto na forma linear (proporcional à velocidade) quanto quadrática (proporcional ao quadrado da velocidade), permitindo avaliar o efeito desse fator sobre o movimento e a acurácia do método numérico.
\end{enumerate}

Para cada caso, os parâmetros de entrada foram lidos dos arquivos de input e utilizados para configurar as simulações. O método de Euler foi empregado para resolver numericamente as equações diferenciais do movimento, permitindo o cálculo iterativo da posição, velocidade e tempo até o corpo atingir o solo.

As simulações foram realizadas para diferentes valores de $\Delta t$ e, nos casos com resistência do ar, para diferentes valores de $c$. Foram analisados os resultados de posições, velocidades e tempos finais obtidos em cada experimento.
\section{Parâmetros de Entrada das Simulações}
\label{sec_parametros_entrada}


\textbf{Tabela 1 — Parâmetros para Queda Livre Simples}

A Tabela 1 apresenta os parâmetros utilizados para a simulação da queda livre simples. Os experimentos foram realizados variando a altura inicial, a velocidade inicial, a gravidade e o passo de tempo, conforme os valores listados.

\begin{table}[H]
\centering
\caption{Tabela – Dados referentes a Entrada queda livre.}
\label{tab:entrada_queda_livre}
\begin{tabular}{llll}
\toprule
H (m) & V (m/s) & g (m/s²) & dt (s) \\
\midrule
100.0 & 0.0 & 9.81 & 1.0 \\
100.0 & 0.0 & 9.81 & 0.5 \\
100.0 & 0.0 & 9.81 & 0.1 \\
100.0 & 0.0 & 9.81 & 0.01 \\
100.0 & 0.0 & 9.81 & 0.001 \\
\bottomrule
\end{tabular}

\end{table}


\textbf{Tabela 2 — Parâmetros para Queda Livre com Resistência Quadrática}

A Tabela 2 mostra os parâmetros empregados na simulação da queda livre considerando resistência do ar quadrática. Foram utilizados diferentes valores para massa, coeficiente de arrasto, área, densidade do ar, tempo máximo e passo de tempo.

\begin{tabular}{llllllll}
\toprule
Modelo & Altura (m) & Massa (kg) & Cd & A (m²) & Rho (kg/m³) & t\_max (s) & dt (s) \\
\midrule
resistencia\_quadratica & 100 & 1.0 & 0.47 & 0.01 & 1.225 & 20 & 1.0 \\
resistencia\_quadratica & 100 & 1.0 & 0.47 & 0.01 & 1.225 & 20 & 0.5 \\
resistencia\_quadratica & 100 & 1.0 & 0.47 & 0.01 & 1.225 & 20 & 0.1 \\
resistencia\_quadratica & 100 & 1.0 & 0.47 & 0.01 & 1.225 & 20 & 0.01 \\
resistencia\_quadratica & 100 & 1.0 & 0.47 & 0.01 & 1.225 & 20 & 0.001 \\
\bottomrule
\end{tabular}


\textbf{Tabela 3 — Parâmetros para Queda Livre com Resistência Linear}

A Tabela 3 apresenta os parâmetros adotados para a simulação da queda livre com resistência do ar linear. Os experimentos consideraram diferentes valores para o coeficiente de resistência, além das demais variáveis relevantes.

\begin{tabular}{llllll}
\toprule
modelo & H (m) & V (m/s) & g (m/s²) & dt (s) & c (1/s) \\
\midrule
resistencia\_linear & 100 & 0.0 & 9.81 & 1.0 & 0.01 \\
resistencia\_linear & 100 & 0.0 & 9.81 & 0.5 & 0.01 \\
resistencia\_linear & 100 & 0.0 & 9.81 & 0.1 & 0.01 \\
resistencia\_linear & 100 & 0.0 & 9.81 & 0.01 & 0.01 \\
resistencia\_linear & 100 & 0.0 & 9.81 & 0.001 & 0.01 \\
resistencia\_linear & 100 & 0.0 & 9.81 & 1.0 & 0.1 \\
resistencia\_linear & 100 & 0.0 & 9.81 & 0.5 & 0.1 \\
resistencia\_linear & 100 & 0.0 & 9.81 & 0.1 & 0.1 \\
resistencia\_linear & 100 & 0.0 & 9.81 & 0.01 & 0.1 \\
resistencia\_linear & 100 & 0.0 & 9.81 & 0.001 & 0.1 \\
resistencia\_linear & 100 & 0.0 & 9.81 & 1.0 & 2.0 \\
resistencia\_linear & 100 & 0.0 & 9.81 & 0.5 & 2.0 \\
resistencia\_linear & 100 & 0.0 & 9.81 & 0.1 & 2.0 \\
resistencia\_linear & 100 & 0.0 & 9.81 & 0.01 & 2.0 \\
resistencia\_linear & 100 & 0.0 & 9.81 & 0.001 & 2.0 \\
resistencia\_linear & 100 & 0.0 & 9.81 & 1.0 & 10.0 \\
resistencia\_linear & 100 & 0.0 & 9.81 & 0.5 & 10.0 \\
resistencia\_linear & 100 & 0.0 & 9.81 & 0.1 & 10.0 \\
resistencia\_linear & 100 & 0.0 & 9.81 & 0.01 & 10.0 \\
resistencia\_linear & 100 & 0.0 & 9.81 & 0.001 & 10.0 \\
\bottomrule
\end{tabular}

\section{Observação dos Resultados, Cálculos Analíticos, Análise de Acurácia e Estatística}
\label{sec_analise_resultados}

Esta seção reúne toda a análise dos resultados, os cálculos analíticos, a avaliação da acurácia dos métodos numéricos utilizados e a análise estatística dos experimentos. A análise foi realizada a partir da comparação entre os valores numéricos obtidos nas simulações e as soluções analíticas conhecidas para cada caso físico. As principais equações utilizadas foram:

\begin{itemize}
    \item \textbf{Velocidade final sem resistência do ar:}
    \begin{equation}
        v_{\text{final}} = V - g \cdot t
    \end{equation}
    onde $V$ é a velocidade inicial, $g$ a aceleração da gravidade e $t$ o tempo total de queda.
    \item \textbf{Velocidade final com resistência do ar linear:}
    \begin{equation}
        v_{\text{final}} = (V + \frac{g}{c}) \cdot e^{c t} - \frac{g}{c}
    \end{equation}
    onde $c$ é o coeficiente de resistência do ar (linear).
    \item \textbf{Erro relativo:}
    \begin{equation}
        E = \left| \frac{v_{\text{num}} - v_{\text{ana}}}{v_{\text{ana}}} \right|
    \end{equation}
    onde $E$ representa o erro relativo, $v_{\text{num}}$ é o valor numérico obtido na simulação e $v_{\text{ana}}$ o valor analítico.
\end{itemize}

Essas fórmulas foram implementadas em funções Python para automatizar o cálculo da acurácia dos métodos numéricos. A acurácia foi avaliada pela variação do erro relativo médio em função do passo de tempo ($\Delta t$) e do coeficiente de resistência do ar, permitindo identificar a convergência e a precisão do método empregado. Foram gerados gráficos dos principais resultados, incluindo:
\begin{itemize}
    \item Posição em função do tempo;
    \item Velocidade em função do tempo;
    \item Aceleração em função do tempo;
    \item Erro relativo em função do passo de tempo ($\Delta t$).
\end{itemize}
Esses gráficos permitiram visualizar o comportamento do sistema sob diferentes condições e validar os modelos computacionais empregados.

A análise estatística dos resultados envolveu a comparação entre as soluções numéricas obtidas pelo método de Euler e a solução analítica do problema da queda livre. Foram avaliados diferentes valores de passo de tempo ($\Delta t$) para investigar o impacto na precisão dos resultados. O erro absoluto entre o tempo de queda simulado e o valor analítico foi calculado e representado graficamente, permitindo avaliar a acurácia do método numérico. Além disso, foram realizadas simulações considerando a resistência do ar, possibilitando a comparação entre diferentes cenários físicos e a análise do comportamento do sistema sob diferentes hipóteses.
